\begin{abstract}

\begin{sloppy}
Энэхүү дипломын ажлын зорилго нь Монголын аялал жуулчлалын туршлагыг илүү оролцоонд суурилсан, интерактив, танин мэдэхүйн үнэ цэнтэй хэлбэрт шилжүүлэх Unity төвтэй AR тоглоомын прототип системийг зохиомжлон, хөгжүүлж, үнэлэхэд оршино. Судалгаанд өргөтгөсөн бодит байдал (AR), байршилд суурилсан үйлчилгээ (LBS), тоглоомжуулалт, хэрэглэгчийн оролцооны онол, мөн мэдрэгчийн нэгтгэл ба орон зайн байршуулалтын техникийн суурийг нэгтгэн авч үзсэн.

Системийн архитектур нь Unity клиент, Node.js/Express backend, PostgreSQL/Prisma өгөгдлийн сан гэсэн үндсэн бүрэлдэхүүнтэй. Хэрэглэгч байршлын радиуст ормогц AR туршлага идэвхжиж, виртуал авдар нээх замаар байршилтай уялдсан reward болон соёлын тайлбар хүлээн авах боломжтой. Зохиомжийн үндсэн инженерийн санаа нь GPS-д суурилсан макро түвшний идэвхжүүлэлт ба AR tracking-д суурилсан микро түвшний динамик байршуулалтыг салгаж үзсэнд оршино. Энэ шийдэл нь GPS-ийн хэлбэлзлээс үүдэх UX-ийн доголдлыг бууруулж, виртуал объектыг хэрэглэгчид илүү ойлгомжтой байдлаар дүрслэх боломж бүрдүүлэсэн.
\end{sloppy}

\end{abstract}
